\section{Limitations of this analysis}

The analysis has four main limitations.

First, the description of the workflow rests on documents and on one source of practice.
It draws on the model's documentation, code and 2023 data, and on a recorded working session with one of the developers.
How long the steps take and where the judgment lies come largely from that account, and were not observed across several analysts or studies.
Other analysts may divide the work differently.

Second, the evidence on agents is young.
Most of the studies in Chapters 2 and 4 are recent preprints, many describe a single system evaluated by its own authors, and results on real systems are weaker than on synthetic ones (Section 2.6).
The claims of this thesis about what an agent could do in each step are inferences from that evidence, not measurements on GGM.
Language models also change quickly, so statements about their capabilities may date within a year.
The division of labour argued for here depends less on them: it rests on which tasks have a checkable answer and which carry economic judgment, and that does not change with a better model.

Third, the analysis concerns one model.
Many energy and economic models share GGM's features, namely large spreadsheet inputs, a calibration to reference values, and parameters that encode a view of the market, and the method of Chapter 5 should transfer to them.
That it does has been argued, not shown.

Fourth, nothing has been built or tested.
The system in Chapter 5 is a design, and its costs, the time it would save and the quality of its suggestions are unknown until it is built.
This was the intended scope of the project thesis, and the tests in Table \ref{tab:agents} are designed to settle these questions in the master's thesis.
